\newpage
\thispagestyle{empty}
\mbox{}
\newpage

\addchapnonumber{B.\ \ What research in 2050?}
\label{sec:research_futur}
% recall the intro about global warming
This appendix chapter follows the introduction section about global warming (\cref{sec:anthropo_disturbance}}), and more precisely, adresses the question risen in \cref{sec:implication_sciences} about the implications of global warming near future consequences for the scientific research.
The objective of this appendix is to express some of the elements that have driven my thoughts during the three years of my PhD. 
They take their origins in discussions with young and less young colleagues, in seminars I followed, in meetings with concerned and engaged researchers (from my university and beyond), and most of all in the wavy tides of feelings that I had all along my PhD years and before.
\vspace{\baselineskip}

Even though social sciences rarely link political events to climate conditions at first sight, historical reviews enlighten by climate sciences reveal a strong influence of climate changes on human societies. The French historian Jean-Baptiste Fressoz gives evidences that climate was a key actor in every major historical event \citep{fressoz_revoltes_2020}.
Similarly, the history since the industrial revolutions has been driven by fossil fuels. In his book \textit{Or noir} \citep{auzanneau_or_2015}, Matthieu Auzanneau demonstrates the central role of petrol on human recent history, from the social organisation and economical management of natural and human resources, to the outcome of World Wars (WWs).
There is no scientifically demonstrated reasons for this to change in the future. The recent international geopolitical news show it.
% ce que ça va casser

\vspace{\baselineskip}
\textbf{Resilience of the research system}\\
In 2024, two events threatened the researchs in Earth observations from space and planetology. 
The first one was the proposition for a moratorium on new space missions within the IPSL\footnote{The IPSL gathers the biggest Earth sciences laboratories of the Parisian region.} laboratories.
This suggestion issued from a reflexion group made of IPSL researchers, charged to find levers that would allow the drastic reduction of the environmental footprint of our research activities.
Indeed, in a IPSL laboratory, the carbon footprint share of a researcher is between \num{10} and \SI{15}{\tonne\ce{CO2}eq\per\year}. 
It accounts only for the professionnal part, and yet it is at least five times the goal of \SI{2}{\tonne\ce{CO2}eq\per\year}.
The same calculation at the scale of CNRS, all fields combined, gives the same order of magnitude.
%so yes, no more space mission was not stupid
So, the idea was not pointless. The carbon footprint of space sciences has been evaluated by \citet{knodlseder_estimate_2022}, and appears to be the first emission category. In LATMOS for example, the carbon footprints of solar system and Earth observation space missions reachs the same value as all the other activities, i.e. around \tcoeq{10} per capita and per year.
% but people said it was, and that it was impossible, and that it would kill the science, and the competition, etc.
But people reacted as it was senseless, since it would be impossible to stop the work of the concerned people, and that the evolution of their research practices could not be done.
The argument was efficient, and it made the proposition almost unanimously rejected by the community.

    % the second threat was Trump
The second event that happened in 2024, was the election of D. Trump as president of the U.S.A. in December. 
Its political line led to cancel tenth of Nasa's and NOAA's programs. It impacted researchers in a great range of scientific fields.
Some cancellation appeared to be even more impactful that what would have been the moratorium on new space missions.
Nevertheless, the concerned communities organised in order to face this new reality collectively. And yet, research goes on; differently that is for sure, but events of 2025 have shown that it is possible to change our research practices.

It have also demonstrated that it is preferable to do it on purpose, rather than suffer it.
Because poverty and sufficiency are the same, the only difference is that the first is not chosen when the second is, and it makes all the difference.
% he do not like fundamental sciences so he cut budgets
% the community adapted

% but in the future, it is likely to get more of these situations

\vspace{\baselineskip}
\textbf{Researchers commitment}\\
In climate sciences, researchers have access\---through their skills and knowledge\---to the understanding of how the world is and will be. They have the privilege of getting, by means of their research, a broader look to our overwhelming modern daily life. In other words, we have access to the meaning of the IPCC report conclusions. 
The question thus is, is it of our responsibility to take action consequently? And, if not us, who else?

    % That said, what kind of research do we fancy for the future?
Once this question is asked, different pathways of action open up. One of them is the civic commitment. Researchers are citizens who could want to get engaged and to put their knowledge to good use, but ``how to do it'' remains an issue.
About climate change and life conditions collapse, some colleagues get in touch with \textit{Scientist Rebellion}, and take action with them. Others write books, teach, talk on TV, or ride their bike from town to town in order to tell the climate to young children.
There is no best behaviour, but people will regularly ask themselves if they could do better or more.

In order to investigate this interrogation, one could do the following experiment:
\emph{Let's put ourselves in a 20 years old person's shoes, in 2070. We live in an occidental Europe that have endured the consequences of almost three degrees of global warming. 
\tempc{50} is now common in southern France, and no Spanish forest are left to burn. The agricultural production has become a major concern since the natural water supplies became chaotic.
The European governments have folded in on themselves and became nationalist while facing the constantly increasing flow of war and climate refugees.
The economic growth is long gone, as predicted one century ago by the Meadows. Only remains agriculture and war industry on the international market.} 

\emph{We are this young person who, for the first time in many generations of their family, will not have the possibility to pursue studies, and maybe not even to go watch a film on a Friday night. % je ne suis pas tellement convaincu de ça finalement
This young European thinks about the past, about the people in 2020 and their actions. Do they knew what would have been the future? Were they able to foresee it? What did they do to prevent the worse? 
About those very engaged people, those among the most knowledgeable: were their actions commensurate with what was at stake? In their place, what would I have done?}


\vspace{\baselineskip}
Another pathway of action is to think about the evolution of research. For sure it will greatly suffer of the future climate conditions, and will need to adapt in order to maintain itself.
In twenty years, the world would have change so much, but will we have anticipated anything or just endured everything?
Is not it our responsability, as being the most climate knowledgeable citizens, to use our understanding of the present in order to fancy the future, and therefore take action to change\---if not the entire society\---our communities of researchers?
Because, if we cannot face successfully, for ourselves, the challenges of the 21st century, who will?


    % does research make decreasing the entropy?

% --> opening question: what the research will be like in 2050?
